%!TeX root=../emmatop.tex
\chapter[Chapter \thechapter]{}
\lettrine[lraise=0.3]{E}{mma} and Harriet had been walking together one morning, and, in Emma's opinion, had been talking enough of Mr Elton for that day. She could not think that Harriet's solace or her own sins required more; and she was therefore industriously getting rid of the subject as they returned;—but it burst out again when she thought she had succeeded, and after speaking some time of what the poor must suffer in winter, and receiving no other answer than a very plaintive—<Mr Elton is so good to the poor!> she found something else must be done.

They were just approaching the house where lived Mrs and Miss Bates. She determined to call upon them and seek safety in numbers. There was always sufficient reason for such an attention; Mrs and Miss Bates loved to be called on, and she knew she was considered by the very few who presumed ever to see imperfection in her, as rather negligent in that respect, and as not contributing what she ought to the stock of their scanty comforts.

She had had many a hint from Mr Knightley and some from her own heart, as to her deficiency—but none were equal to counteract the persuasion of its being very disagreeable,—a waste of time—tiresome women—and all the horror of being in danger of falling in with the second-rate and third-rate of Highbury, who were calling on them for ever, and therefore she seldom went near them. But now she made the sudden resolution of not passing their door without going in—observing, as she proposed it to Harriet, that, as well as she could calculate, they were just now quite safe from any letter from Jane Fairfax.

The house belonged to people in business. Mrs and Miss Bates occupied the drawing-room floor; and there, in the very moderate-sized apartment, which was every thing to them, the visitors were most cordially and even gratefully welcomed; the quiet neat old lady, who with her knitting was seated in the warmest corner, wanting even to give up her place to Miss Woodhouse, and her more active, talking daughter, almost ready to overpower them with care and kindness, thanks for their visit, solicitude for their shoes, anxious inquiries after Mr Woodhouse's health, cheerful communications about her mother's, and sweet-cake from the beaufet—<Mrs Cole had just been there, just called in for ten minutes, and had been so good as to sit an hour with them, and she had taken a piece of cake and been so kind as to say she liked it very much; and, therefore, she hoped Miss Woodhouse and Miss Smith would do them the favour to eat a piece too.>

The mention of the Coles was sure to be followed by that of Mr Elton. There was intimacy between them, and Mr Cole had heard from Mr Elton since his going away. Emma knew what was coming; they must have the letter over again, and settle how long he had been gone, and how much he was engaged in company, and what a favourite he was wherever he went, and how full the Master of the Ceremonies' ball had been; and she went through it very well, with all the interest and all the commendation that could be requisite, and always putting forward to prevent Harriet's being obliged to say a word.

This she had been prepared for when she entered the house; but meant, having once talked him handsomely over, to be no farther incommoded by any troublesome topic, and to wander at large amongst all the Mistresses and Misses of Highbury, and their card-parties. She had not been prepared to have Jane Fairfax succeed Mr Elton; but he was actually hurried off by Miss Bates, she jumped away from him at last abruptly to the Coles, to usher in a letter from her niece.

<Oh! yes—Mr Elton, I understand—certainly as to dancing—Mrs Cole was telling me that dancing at the rooms at Bath was—Mrs Cole was so kind as to sit some time with us, talking of Jane; for as soon as she came in, she began inquiring after her, Jane is so very great a favourite there. Whenever she is with us, Mrs Cole does not know how to shew her kindness enough; and I must say that Jane deserves it as much as any body can. And so she began inquiring after her directly, saying, <I know you cannot have heard from Jane lately, because it is not her time for writing;> and when I immediately said, <But indeed we have, we had a letter this very morning,> I do not know that I ever saw any body more surprized. <Have you, upon your honour?> said she; <well, that is quite unexpected. Do let me hear what she says.>>

Emma's politeness was at hand directly, to say, with smiling interest—

<Have you heard from Miss Fairfax so lately? I am extremely happy. I hope she is well?>

\begin{figure}[tbph]
\centering
\includegraphics[width=\linewidth]{19hereitis}
\caption{<Oh! here it is.>}
\end{figure}

<Thank you. You are so kind!> replied the happily deceived aunt, while eagerly hunting for the letter.—<Oh! here it is. I was sure it could not be far off; but I had put my huswife upon it, you see, without being aware, and so it was quite hid, but I had it in my hand so very lately that I was almost sure it must be on the table. I was reading it to Mrs Cole, and since she went away, I was reading it again to my mother, for it is such a pleasure to her—a letter from Jane—that she can never hear it often enough; so I knew it could not be far off, and here it is, only just under my huswife—and since you are so kind as to wish to hear what she says;—but, first of all, I really must, in justice to Jane, apologise for her writing so short a letter—only two pages you see—hardly two—and in general she fills the whole paper and crosses half. My mother often wonders that I can make it out so well. She often says, when the letter is first opened, <Well, Hetty, now I think you will be put to it to make out all that checker-work>—don't you, ma'am?—And then I tell her, I am sure she would contrive to make it out herself, if she had nobody to do it for her—every word of it—I am sure she would pore over it till she had made out every word. And, indeed, though my mother's eyes are not so good as they were, she can see amazingly well still, thank God! with the help of spectacles. It is such a blessing! My mother's are really very good indeed. Jane often says, when she is here, <I am sure, grandmama, you must have had very strong eyes to see as you do—and so much fine work as you have done too!—I only wish my eyes may last me as well.>>

All this spoken extremely fast obliged Miss Bates to stop for breath; and Emma said something very civil about the excellence of Miss Fairfax's handwriting.

<You are extremely kind,> replied Miss Bates, highly gratified; <you who are such a judge, and write so beautifully yourself. I am sure there is nobody's praise that could give us so much pleasure as Miss Woodhouse's. My mother does not hear; she is a little deaf you know. Ma'am,> addressing her, <do you hear what Miss Woodhouse is so obliging to say about Jane's handwriting?>

And Emma had the advantage of hearing her own silly compliment repeated twice over before the good old lady could comprehend it. She was pondering, in the meanwhile, upon the possibility, without seeming very rude, of making her escape from Jane Fairfax's letter, and had almost resolved on hurrying away directly under some slight excuse, when Miss Bates turned to her again and seized her attention.

<My mother's deafness is very trifling you see—just nothing at all. By only raising my voice, and saying any thing two or three times over, she is sure to hear; but then she is used to my voice. But it is very remarkable that she should always hear Jane better than she does me. Jane speaks so distinct! However, she will not find her grandmama at all deafer than she was two years ago; which is saying a great deal at my mother's time of life—and it really is full two years, you know, since she was here. We never were so long without seeing her before, and as I was telling Mrs Cole, we shall hardly know how to make enough of her now.>

<Are you expecting Miss Fairfax here soon?>

<Oh yes; next week.>

<Indeed!—that must be a very great pleasure.>

<Thank you. You are very kind. Yes, next week. Every body is so surprized; and every body says the same obliging things. I am sure she will be as happy to see her friends at Highbury, as they can be to see her. Yes, Friday or Saturday; she cannot say which, because Colonel Campbell will be wanting the carriage himself one of those days. So very good of them to send her the whole way! But they always do, you know. Oh yes, Friday or Saturday next. That is what she writes about. That is the reason of her writing out of rule, as we call it; for, in the common course, we should not have heard from her before next Tuesday or Wednesday.>

<Yes, so I imagined. I was afraid there could be little chance of my hearing any thing of Miss Fairfax to-day.>

<So obliging of you! No, we should not have heard, if it had not been for this particular circumstance, of her being to come here so soon. My mother is so delighted!—for she is to be three months with us at least. Three months, she says so, positively, as I am going to have the pleasure of reading to you. The case is, you see, that the Campbells are going to Ireland. Mrs Dixon has persuaded her father and mother to come over and see her directly. They had not intended to go over till the summer, but she is so impatient to see them again—for till she married, last October, she was never away from them so much as a week, which must make it very strange to be in different kingdoms, I was going to say, but however different countries, and so she wrote a very urgent letter to her mother—or her father, I declare I do not know which it was, but we shall see presently in Jane's letter—wrote in Mr Dixon's name as well as her own, to press their coming over directly, and they would give them the meeting in Dublin, and take them back to their country seat, Baly-craig, a beautiful place, I fancy. Jane has heard a great deal of its beauty; from Mr Dixon, I mean—I do not know that she ever heard about it from any body else; but it was very natural, you know, that he should like to speak of his own place while he was paying his addresses—and as Jane used to be very often walking out with them—for Colonel and Mrs Campbell were very particular about their daughter's not walking out often with only Mr Dixon, for which I do not at all blame them; of course she heard every thing he might be telling Miss Campbell about his own home in Ireland; and I think she wrote us word that he had shewn them some drawings of the place, views that he had taken himself. He is a most amiable, charming young man, I believe. Jane was quite longing to go to Ireland, from his account of things.>

At this moment, an ingenious and animating suspicion entering Emma's brain with regard to Jane Fairfax, this charming Mr Dixon, and the not going to Ireland, she said, with the insidious design of farther discovery,

<You must feel it very fortunate that Miss Fairfax should be allowed to come to you at such a time. Considering the very particular friendship between her and Mrs Dixon, you could hardly have expected her to be excused from accompanying Colonel and Mrs Campbell.>

<Very true, very true, indeed. The very thing that we have always been rather afraid of; for we should not have liked to have her at such a distance from us, for months together—not able to come if any thing was to happen. But you see, every thing turns out for the best. They want her (Mr and Mrs Dixon) excessively to come over with Colonel and Mrs Campbell; quite depend upon it; nothing can be more kind or pressing than their joint invitation, Jane says, as you will hear presently; Mr Dixon does not seem in the least backward in any attention. He is a most charming young man. Ever since the service he rendered Jane at Weymouth, when they were out in that party on the water, and she, by the sudden whirling round of something or other among the sails, would have been dashed into the sea at once, and actually was all but gone, if he had not, with the greatest presence of mind, caught hold of her habit— (I can never think of it without trembling!)—But ever since we had the history of that day, I have been so fond of Mr Dixon!>

<But, in spite of all her friends' urgency, and her own wish of seeing Ireland, Miss Fairfax prefers devoting the time to you and Mrs Bates?>

<Yes—entirely her own doing, entirely her own choice; and Colonel and Mrs Campbell think she does quite right, just what they should recommend; and indeed they particularly wish her to try her native air, as she has not been quite so well as usual lately.>

<I am concerned to hear of it. I think they judge wisely. But Mrs Dixon must be very much disappointed. Mrs Dixon, I understand, has no remarkable degree of personal beauty; is not, by any means, to be compared with Miss Fairfax.>

<Oh! no. You are very obliging to say such things—but certainly not. There is no comparison between them. Miss Campbell always was absolutely plain—but extremely elegant and amiable.>

<Yes, that of course.>

<Jane caught a bad cold, poor thing! so long ago as the 7th of November, (as I am going to read to you,) and has never been well since. A long time, is not it, for a cold to hang upon her? She never mentioned it before, because she would not alarm us. Just like her! so considerate!—But however, she is so far from well, that her kind friends the Campbells think she had better come home, and try an air that always agrees with her; and they have no doubt that three or four months at Highbury will entirely cure her—and it is certainly a great deal better that she should come here, than go to Ireland, if she is unwell. Nobody could nurse her, as we should do.>

<It appears to me the most desirable arrangement in the world.>

<And so she is to come to us next Friday or Saturday, and the Campbells leave town in their way to Holyhead the Monday following—as you will find from Jane's letter. So sudden!—You may guess, dear Miss Woodhouse, what a flurry it has thrown me in! If it was not for the drawback of her illness—but I am afraid we must expect to see her grown thin, and looking very poorly. I must tell you what an unlucky thing happened to me, as to that. I always make a point of reading Jane's letters through to myself first, before I read them aloud to my mother, you know, for fear of there being any thing in them to distress her. Jane desired me to do it, so I always do: and so I began to-day with my usual caution; but no sooner did I come to the mention of her being unwell, than I burst out, quite frightened, with <Bless me! poor Jane is ill!>—which my mother, being on the watch, heard distinctly, and was sadly alarmed at. However, when I read on, I found it was not near so bad as I had fancied at first; and I make so light of it now to her, that she does not think much about it. But I cannot imagine how I could be so off my guard. If Jane does not get well soon, we will call in Mr Perry. The expense shall not be thought of; and though he is so liberal, and so fond of Jane that I dare say he would not mean to charge any thing for attendance, we could not suffer it to be so, you know. He has a wife and family to maintain, and is not to be giving away his time. Well, now I have just given you a hint of what Jane writes about, we will turn to her letter, and I am sure she tells her own story a great deal better than I can tell it for her.>

<I am afraid we must be running away,> said Emma, glancing at Harriet, and beginning to rise—<My father will be expecting us. I had no intention, I thought I had no power of staying more than five minutes, when I first entered the house. I merely called, because I would not pass the door without inquiring after Mrs Bates; but I have been so pleasantly detained! Now, however, we must wish you and Mrs Bates good morning.>

And not all that could be urged to detain her succeeded. She regained the street—happy in this, that though much had been forced on her against her will, though she had in fact heard the whole substance of Jane Fairfax's letter, she had been able to escape the letter itself.